% \documentclass[tikz,border=10pt]{standalone}
% \usepackage{amsmath}
% \usepackage{tikz}
% \usetikzlibrary{angles, quotes, calc}

% \begin{document}
% \begin{tikzpicture}[scale=1]

% % Vektorlänge und Abstand
% \def\veclenA{2.5}
% \def\veclenB{1.75}
% \def\spacing{4}
% \def\angleradius{20}

% % Winkeldefinitionen
% \def\angleA{120}
% \def\angleB{90}
% \def\angleC{45}
% \def\angleD{0}

% % Farben (Okabe Ito Palette, https://f0nzie.github.io/dataviz-rsuite/color-pitfalls.html)
% \definecolor{veca}{HTML}{0072B2}  % Blau für a
% % \definecolor{vecb}{HTML}{E69F00}  % Senfgelb für b
% \definecolor{vecb}{HTML}{D55E00}  % kräftigeres, rötlicheres Orange für b
% % \definecolor{proj}{HTML}{009E73}  % bluish green für Projektion
% \definecolor{proj}{HTML}{CC79A7}  % reddish purple für Projektion
% \definecolor{anglec}{HTML}{009E73} % bluish green für Winkel
% % \definecolor{anglec}{HTML}{66C2A5} % leicht entsättigtes Türkis für Winkel
% % \definecolor{anglec}{HTML}{CC79A7} % reddish purple für Winkel
% % \definecolor{helper}{HTML}{BBBBBB} % Hellgrau für Hilfslinien
% \definecolor{helper}{HTML}{E69F00} % orange für Hilfslinien

% % Szene 1: alpha = 120°
% \def\angle{\angleA}
% \begin{scope}[shift={(0,0)}]
%   \coordinate (O) at (0,0);
%   \coordinate (A) at (\veclenA,0);
%   \coordinate (B) at (\angle:\veclenB);
%   \coordinate (P) at ({\veclenB*cos(\angle)},0);

%   \node at ($(O)!0.5!(A) + (0,\veclenA)$) {$\alpha=\angle^\circ$};
%   \draw[->, thick, veca] (O) -- (A) node[right] {$\vec{a}$};
%   \draw[->, thick, vecb] (O) -- (B) node[above] {$\vec{b}$};
%   \draw[dashed, helper] (B) -- (P);
%   \draw[thick, proj] (P) -- (O);
%   \node at ($(O)!0.5!(P) -(0,0.2*\veclenA)$) {$\vec{a} \cdot \vec{b} < 0$};
%   \pic[draw, color=anglec, ->, angle radius=\angleradius, "$\alpha$"] {angle=A--O--B};
% \end{scope}

% % Szene 2: alpha = 90°
% \def\angle{\angleB}
% \begin{scope}[shift={(\spacing,0)}]
%   \coordinate (O) at (0,0);
%   \coordinate (A) at (\veclenA,0);
%   \coordinate (B) at (\angle:\veclenB);
%   \coordinate (P) at ({\veclenB*cos(\angle)},0);

%   \node at ($(O)!0.5!(A) + (0,\veclenA)$) {$\alpha=\angle^\circ$};
%   \draw[->, thick, veca] (O) -- (A) node[right] {$\vec{a}$};
%   \draw[->, thick, vecb] (O) -- (B) node[above] {$\vec{b}$};
%   \draw[dashed, helper] (B) -- (P);
%   % \draw[thick, proj] (P) -- (O);
%   \fill[proj] (O) circle (1.5pt);
%   \node at ($(O)!0.5!(P) -(0,0.2*\veclenA)$) {$\vec{a} \cdot \vec{b} = 0$};
%   \pic[draw, color=anglec, ->, angle radius=\angleradius, "$\alpha$"] {angle=A--O--B};
% \end{scope}

% % Szene 3: alpha = 45°
% \def\angle{\angleC}
% \begin{scope}[shift={(2*\spacing,0)}]
%   \coordinate (O) at (0,0);
%   \coordinate (A) at (\veclenA,0);
%   \coordinate (B) at (\angle:\veclenB);
%   \coordinate (P) at ({\veclenB*cos(\angle)},0);

%   \node at ($(O)!0.5!(A) + (0,\veclenA)$) {$\alpha=\angle^\circ$};
%   \draw[->, thick, veca] (O) -- (A) node[right] {$\vec{a}$};
%   \draw[->, thick, vecb] (O) -- (B) node[above] {$\vec{b}$};
%   \draw[dashed, helper] (B) -- (P);
%   \draw[thick, proj] (P) -- (O);
%   \node at ($(O)!0.5!(P) -(0,0.2*\veclenA)$) {$\vec{a} \cdot \vec{b} > 0$};
%   \pic[draw, color=anglec, ->, angle radius=\angleradius, "$\alpha$"] {angle=A--O--B};
% \end{scope}

% % Szene 4: alpha = 0°
% \def\angle{\angleD}
% \begin{scope}[shift={(3*\spacing,0)}]
%   \coordinate (O) at (0,0);
%   \coordinate (A) at (\veclenA,0);
%   \coordinate (B) at (\angle:\veclenB);
%   \coordinate (P) at ({\veclenB*cos(\angle)},0);

%   \node at ($(O)!0.5!(A) + (0,\veclenA)$) {$\alpha=\angle^\circ$};
%   \draw[->, thick, veca] (O) -- (A) node[right] {$\vec{a}$};
%   \draw[->, thick, vecb] (O) -- (B) node[above] {$\vec{b}$};
%   % \draw[dashed, helper] (B) -- (P);
%   \draw[thick, proj] (P) -- (O);
%   \node at ($(O)!0.5!(P) -(0,0.2*\veclenA)$) {$\vec{a} \cdot \vec{b} = |\vec{a}| |\vec{b}|$};
%   % \pic[draw, color=anglec, ->, angle radius=\angleradius, "$\alpha$"] {angle=A--O--B};
% \end{scope}

% \end{tikzpicture}
% \end{document}





\documentclass[tikz,border=10pt]{standalone}
\usepackage{amsmath}
\usepackage{tikz}
\usetikzlibrary{angles, quotes, calc}

\begin{document}
\begin{tikzpicture}[scale=1]

% Vektorlänge und Abstand
\def\veclenA{2.5}
\def\veclenB{1.75}
\def\spacing{4}
\def\angleradius{20}

% Farben (Okabe Ito Palette, https://f0nzie.github.io/dataviz-rsuite/color-pitfalls.html)
\definecolor{veca}{HTML}{0072B2}  % Blau für a
% \definecolor{vecb}{HTML}{E69F00}  % Senfgelb für b
\definecolor{vecb}{HTML}{D55E00}  % kräftigeres, rötlicheres Orange für b
% \definecolor{proj}{HTML}{009E73}  % bluish green für Projektion
\definecolor{proj}{HTML}{CC79A7}  % reddish purple für Projektion
\definecolor{anglec}{HTML}{009E73} % bluish green für Winkel
% \definecolor{anglec}{HTML}{66C2A5} % leicht entsättigtes Türkis für Winkel
% \definecolor{anglec}{HTML}{CC79A7} % reddish purple für Winkel
% \definecolor{helper}{HTML}{BBBBBB} % Hellgrau für Hilfslinien
\definecolor{helper}{HTML}{E69F00} % orange für Hilfslinien


%%%%%%%%%%%%%%%%%%% alpha = 0°
\def\angle{0}
\begin{scope}[shift={(0,0)}]
  \coordinate (O) at (0,0);
  \coordinate (A) at (\veclenA,0);
  \coordinate (B) at (\angle:\veclenB);
  \coordinate (P) at ({\veclenB*cos(\angle)},0);

  \node at ($(O)!0.5!(A) + (0,1.1*\veclenB)$) {$\alpha=\angle^\circ$};
  \draw[->, thick, veca] (O) -- (A) node[above] {$\vec{a}$};
  \draw[->, thick, vecb] (O) -- (B) node[above] {$\vec{b}$};
  % \pic[draw, color=anglec, ->, angle radius=\angleradius, "$\alpha$"] {angle=A--O--B};
  % \draw[dashed, helper] (B) -- (P);
  \draw[thick, proj] (P) -- (O);
  \fill[proj, opacity=0.3] (O) rectangle (P |- 0,-\veclenA);
  \node at ($(O)!0.5!(P) - (0,0.3*\veclenA)$) {$\vec{a} \cdot \vec{b} = |\vec{a}| |\vec{b}|$};
\end{scope}

%%%%%%%%%%%%%%%%%%% alpha = 45°
\def\angle{45}
\begin{scope}[shift={(\spacing,0)}]
  \coordinate (O) at (0,0);
  \coordinate (A) at (\veclenA,0);
  \coordinate (B) at (\angle:\veclenB);
  \coordinate (P) at ({\veclenB*cos(\angle)},0);

  \node at ($(O)!0.5!(A) + (0,1.1*\veclenB)$) {$\alpha=\angle^\circ$};
  \draw[->, thick, veca] (O) -- (A) node[above left] {$\vec{a}$};
  \draw[->, thick, vecb] (O) -- (B) node[right] {$\vec{b}$};
  \pic[draw, color=anglec, ->, angle radius=\angleradius, "$\alpha$"] {angle=A--O--B};
  \draw[dashed, helper] (B) -- (P);
  \draw[thick, proj] (P) -- (O);
  \fill[proj, opacity=0.3] (O) rectangle (P |- 0,-\veclenA);
  \node at ($(O)!0.5!(P) - (0,0.3*\veclenA)$) {$\vec{a} \cdot \vec{b} > 0$};
\end{scope}

%%%%%%%%%%%%%%%%%%% alpha = 90°
\def\angle{90}
\begin{scope}[shift={(2*\spacing,0)}]
  \coordinate (O) at (0,0);
  \coordinate (A) at (\veclenA,0);
  \coordinate (B) at (\angle:\veclenB);
  \coordinate (P) at ({\veclenB*cos(\angle)},0);

  \node at ($(O)!0.5!(A) + (0,1.1*\veclenB)$) {$\alpha=\angle^\circ$};
  \draw[->, thick, veca] (O) -- (A) node[above left] {$\vec{a}$};
  \draw[->, thick, vecb] (O) -- (B) node[right] {$\vec{b}$};
  \pic[draw, color=anglec, ->, angle radius=\angleradius, "$\alpha$"] {angle=A--O--B};
  \draw[dashed, helper] (B) -- (P);
  % \draw[thick, proj] (P) -- (O);
  \fill[proj] (O) circle (1.5pt);
  \node at ($(O)!0.5!(P) - (0,0.3*\veclenA)$) {$\vec{a} \cdot \vec{b} = 0$};
\end{scope}

%%%%%%%%%%%%%%%%%%% alpha = 120°
\def\angle{120}
\begin{scope}[shift={(3*\spacing,0)}]
  % \coordinate (O) at (0,0);
  % \coordinate (A) at (\veclenA,0);
  % \coordinate (B) at (\angle:\veclenB);
  % \coordinate (P) at ({\veclenB*cos(\angle)},0);
  \coordinate (O) at ({-\veclenB*cos(\angle)},0);
  \coordinate (A) at ({\veclenA - \veclenB*cos(\angle)},0);
  \coordinate (B) at (0,{\veclenB*sin(\angle)});
  \coordinate (P) at (0,0);

  \node at ($(P)!0.5!(A) + (0,1.1*\veclenB)$) {$\alpha=\angle^\circ$};
  \draw[->, thick, veca] (O) -- (A) node[above left] {$\vec{a}$};
  \draw[->, thick, vecb] (O) -- (B) node[above right] {$\vec{b}$};
  \pic[draw, color=anglec, ->, angle radius=\angleradius, "$\alpha$"] {angle=A--O--B};
  \draw[dashed, helper] (B) -- (P);
  \draw[thick, proj] (P) -- (O);
  \fill[proj, opacity=0.3] (O) rectangle (P |- 0,-\veclenA);
  \node at ($(O)!0.5!(P) - (0,0.3*\veclenA)$) {$\vec{a} \cdot \vec{b} < 0$};
\end{scope}

%%%%%%%%%%%%%%%%%%% alpha = 180°
\def\angle{180}
\begin{scope}[shift={(4*\spacing,0)}]
  % \coordinate (O) at (0,0);
  % \coordinate (A) at (\veclenA,0);
  % \coordinate (B) at (\angle:\veclenB);
  % \coordinate (P) at ({\veclenB*cos(\angle)},0);
  \coordinate (O) at ({-\veclenB*cos(\angle)},0);
  \coordinate (A) at ({\veclenA - \veclenB*cos(\angle)},0);
  \coordinate (B) at (0,{\veclenB*sin(\angle)});
  \coordinate (P) at (0,0);

  \node at ($(P)!0.5!(A) + (0,1.1*\veclenB)$) {$\alpha=\angle^\circ$};
  \draw[->, thick, veca] (O) -- (A) node[above left] {$\vec{a}$};
  \draw[->, thick, vecb] (O) -- (B) node[above right] {$\vec{b}$};
  \pic[draw, color=anglec, ->, angle radius=\angleradius, "$\alpha$"] {angle=A--O--B};
  \draw[dashed, helper] (B) -- (P);
  \draw[thick, proj] (P) -- (O);
  \fill[proj, opacity=0.3] (O) rectangle (P |- 0,-\veclenA);
  \node at ($(O)!0.5!(P) - (0,0.3*\veclenA)$) {$\vec{a} \cdot \vec{b} = -|\vec{a}| |\vec{b}|$};
\end{scope}


\end{tikzpicture}
\end{document}
